% !PRINTED none
% Abschnitt-opening (part-title) page: no running-head folio visible on this
% scan. Carries the Abschnitt heading, §.175 heading, and body text.

\begin{center}
\gesperrt{Seventeenth Section.}

\gesperrt{Algebraic Solution of Equations.}

\rule{3cm}{0.4pt}

\S.\,175.

\gesperrt{Reduction of the Group by Pure Equations.}
\end{center}

One of the oldest questions, in connection with which the more
recent algebra has chiefly developed, is that of the so-called
\gesperrt{algebraic solution} of equations, by which one understands a
representation of the roots of an equation by a series of
radicals, or their computation by a finite chain of root
extractions. On this question the theory of groups sheds the
brightest light.

If we first make precise the question at issue, it is evidently
this: whether and how one can extend the field $\Omega$ by
successive adjunction of root quantities (radicals) in such a way
that either all or at least a part of the roots
are contained in the extended field. A root quantity, however,
is one which is not itself contained in $\Omega$, but of which some
integral power is; that is, if $a$ is a quantity in $\Omega$,
it is the root of an equation of the form

\[
y^m - a = 0,
\]

that is, of a \gesperrt{pure equation}.

If an \gesperrt{irreducible} equation is to be algebraically solvable,
or at least to have one or several algebraically representable roots,
then after a successive adjunction of roots of
pure equations in finite number the given equation
must become reducible, since a part of its roots is to be contained
in the extended field. Since the initial group $P$ is transitive,
this group must finally become intransitive, or
%CONT>>
